\documentclass[11pt,letterpaper]{article}
\usepackage[top=1.0in,bottom=1.0in,left=1.25in,right=1.25in]{geometry}
\usepackage{microtype}
\usepackage{parskip}
\usepackage{hyperref}
\hypersetup{colorlinks=true,urlcolor=blue,linkcolor=black,citecolor=black}
\setlength{\parindent}{0pt}
\setlength{\parskip}{8pt}

\begin{document}

\textbf{Manuscript:} PONE-D-26-05900\\
\textbf{Title:} Signal Hierarchical Petri Nets: Formal Semantics of
Hierarchical Regulatory Control of Biological Systems\\
\textbf{Date:} \today\\
\textbf{Author:} Sim\~{a}o Eug\'{e}nio

\bigskip
\hrule
\bigskip

Dear Editor,

I am pleased to submit a revised version of the above manuscript in response
to the three reviewers' reports. I am grateful to the editor and to all three
reviewers for their careful and constructive reading. Their comments have
genuinely improved the manuscript, and I hope the revision addresses each
concern satisfactorily.

\medskip
\textbf{Summary of revisions.}
The revision follows directly from the reviewers' recommendations and can be
grouped into three categories:

\textit{(i) Mechanical corrections (Reviewer~2).}
All formal inconsistencies identified by Reviewer~2 have been resolved:
the notation $F_s$ is now used consistently throughout in place of the earlier
$E_s$; the set of transitions $T$ is explicitly introduced at the start of
Definition~1; the test-arc weight convention $W_t \equiv 1$ is stated within
the $F_t$ bullet rather than being left implicit; the Hierarchical Preemption
theorem now opens with an explicit hypothesis on the PreemptionCheck
predicate; and the Irreversible Basin Partition theorem now states the
uniqueness hypothesis explicitly in its statement rather than only in its proof.
A paragraph clarifying SHPN's relationship to classical hybrid Petri nets has
been added, with citations to both the hybrid PN tradition and the specific
extension SHPN introduces.

\textit{(ii) Narrative clarifications (Reviewer~1).}
The term ``first-principles prediction'' has been replaced throughout with
\textit{topology-driven threshold computation}, which more precisely describes
the manuscript's claim: commitment thresholds are structurally localized to at
most two measurable quantities per signal arc, rather than emerging from a
global nonlinear parameter fit. A dual-graph schematic figure has been added
at the close of the Introduction to make the $G_E$\,/\,$G_S$ architecture
visually apparent before any formal notation appears. The relationship to
C/E~nets, occurrence nets, and process-algebraic approaches has been expanded
in Section~4.4. The Discussion has been revised to state explicitly that
acyclicity applies only to the signal control graph~$G_S$, while $G_E$ remains
fully unrestricted and may contain the feedback cycles common in metabolic and
signalling networks.

\textit{(iii) Second biological system (Reviewer~1).}
The Broader Applicability section has been substantially expanded. The
\textit{lambda} phage lysis-lysogeny switch is analysed in detail, with a
SHPN topology figure showing the RecA\,/\,CI\,/\,Cro commitment cascade. The
\textit{Vibrio fischeri} quorum-sensing circuit and GATA1\,/\,PU.1
haematopoietic commitment are discussed as further examples spanning
prokaryotic and eukaryotic decision-making.

\medskip
\textbf{A note on the balance between Reviewers~1 and~3.}
Reviewer~1 requested a second biological example; Reviewer~3 encouraged
concision. These recommendations pull in opposite directions. I have tried to
honour both: the \textit{lambda} phage analysis is self-contained (figure +
two paragraphs), and the \textit{V.\,fischeri} and GATA1 examples are treated
in a single paragraph each. The net change relative to the submitted version
is approximately six pages, all attributable to the material Reviewer~1
specifically requested. I trust the editor to judge whether this balance is
appropriate; I am prepared to trim further on request.

\medskip
\textbf{On the scope of the formalism.}
I wish to be explicit on one point of framing. SHPN is not presented as a
replacement for, or improvement upon, existing Bio-PN formalisms, process
algebras, or Boolean network approaches. Each of those traditions was designed
for a specific problem and solves it correctly. SHPN addresses a question that
those formalisms were not designed to answer: \textit{where, in the network
topology, does a commitment threshold reside?} The answer --- encoded in the
$G_S$ control graph and the acyclicity theorem --- is an addition to the
modeller's toolkit, not a critique of existing tools. The revision makes this
complementary relationship explicit throughout, particularly in Section~4.4
and the revised Discussion.

\medskip
\textbf{On AI-assisted tools.}
Reviewer~2 raised the question of AI tool usage. In the interest of full
transparency, I used AI-assisted writing tools (GitHub Copilot) as a
collaborative aid for drafting, editing, and formulating arguments. All
scientific content --- the formalism, the theorems and their proofs, the
biological interpretations, and the experimental comparisons --- originates
entirely from my own scientific judgement. The AI tools served to sharpen
prose, not to generate scientific ideas. An explicit disclosure statement has
been added to the Acknowledgments section, consistent with PLOS~ONE policy.

\medskip
\textbf{Conclusion.}
I believe the revised manuscript is substantially stronger in both formal
rigour and biological grounding. I am grateful for the opportunity to revise
and look forward to the editor's decision.

\bigskip

Yours sincerely,

\bigskip\bigskip

Sim\~{a}o Eug\'{e}nio\\
\textit{Universidade Federal de Santa Catarina}

\end{document}
